\documentclass{article}

\usepackage[T1]{fontenc}

\usepackage{booktabs,graphicx,siunitx}

\begin{document}

\section{Granger causality}

\begin{tabular}{llll}
\toprule
Direction & F & df & p \\
\midrule
X$\rightarrow$Y & 122.15 & 4, 111 & $<$.001 \\
Y$\rightarrow$X & 3.42 & 4, 111 & .011 \\
\bottomrule
\end{tabular}

\begin{tabular}{llll}
\toprule
Lag & AIC & BIC & HQ \\
\midrule
1 & $-$5.14 & $-$5.00 & $-$5.08 \\
2 & $-$5.14 & $-$4.90 & $-$5.04 \\
3 & $-$5.17 & $-$4.84 & $-$5.04 \\
4 & $-$5.23 & $-$4.81 & $-$5.06 \\
\bottomrule
\end{tabular}

\includegraphics{figures/01_granger-causality/figure_cross-series.png}

A significant X$\rightarrow$Y p means past values of X help predict Y beyond Y's own past - this is predictive precedence, not proof that X causes Y. Check both directions. The result depends entirely on the lag order (this card uses a max lag of 4): choose it on theory or by an information criterion - vars::VARselect reports the AIC- and BIC-minimising lag - and report it, since different lags can flip the conclusion. Make both series stationary first. Method: R lmtest::grangertest (Granger, 1969). Your significance threshold ($\alpha$) is 0.05.

Granger test X$\rightarrow$Y: F(4,111)=122.15, p $<$ .001; Y$\rightarrow$X: F(4,111)=3.42, p = .011 (lag=4).

Statistical basis: Granger causality F-test (Granger, C. W. J. (1969). "Investigating causal relations by econometric models and cross-spectral methods." Econometrica, 37(3), 424-438.); Implementation (lmtest::grangertest) (Zeileis A, Hothorn T (2002). "Diagnostic Checking in Regression Relationships." R News, 2(3), 7-10.).

\end{document}